\chapter{Amyloidosis}
\index{Amyloidosis}
\label{sec:Amyloidosis}

\section{Overview}


\begin{itemize}[noitemsep]
  \item Amyloidosis is a group of diseases characterized by the extracellular deposition of abnormal, insoluble fibrillar proteins (amyloid) in various tissues and organs, disrupting normal structure and function, with three main types clinically important: AL amyloidosis (immunoglobulin light chain---the most common systemic form, associated with plasma cell dyscrasia/multiple myeloma), AA amyloidosis (secondary/reactive
  \item Serum amyloid A protein derived from chronic inflammation), and ATTR amyloidosis (transthyretin
  \item Hereditary or wild-type/senile systemic amyloidosis). The incidence of AL amyloidosis is 3--5 per million, typically presenting at age 60--70.
\end{itemize}
\section{Causes}
\section{Risk Factors}

\begin{itemize}[noitemsep]
  \item Genetic predisposition and family history
  \item Advancing age
  \item Lifestyle factors (diet, exercise, smoking, alcohol)
  \item Environmental exposures
  \item Underlying medical conditions
  \item Medication use
\end{itemize}
\section{Signs and Symptoms}

\subsection{Common symptoms}
\begin{itemize}[noitemsep]
  \item \item Your symptoms depend on the organs involved: AL amyloidosis commonly presents with nephrotic syndrome (kidney amyloidosis---most common presentation, 70\%), restrictive cardiomyopathy (heart failure with preserved ejection fraction, low-voltage QRS on ECG, thickened ventricular walls on echocardiography with a ``speckled'' appearance), hepatomegaly, peripheral and autonomic neuropathy, macroglossia (tongue enlargement---pathognomonic for AL), periorbital purple spots on the skin, and carpal tunnel syndrome. Pain or discomfort in the affected area    \item Pain or discomfort in the affected area
  \end{itemize}

\subsection{Emergency warning signs}
\begin{itemize}[noitemsep]
  \item Severe or worsening symptoms of amyloidosis require immediate medical evaluation
\end{itemize}
\section{Progression and Stages}
The condition may progress through different stages of severity over time. Early stages often involve mild or subtle symptoms that may be overlooked. Moderate stages involve more noticeable symptoms affecting daily function. Advanced stages may involve significant impairment and complications.

\begin{itemize}[noitemsep]
  \item Treatment for AL amyloidosis is anti-plasma cell therapy (similar to multiple myeloma: bortezomib-based regimens, autologous stem cell transplantation in selected patients)
  \item AA amyloidosis is treated by controlling the underlying inflammatory condition
  \item ATTR amyloidosis doctors typically treat it with tafamidis (stabilizes transthyretin, slows neuropathy and cardiomyopathy progression).
\end{itemize}
\section{Complications}
\begin{itemize}[noitemsep]
  \item Disease progression leading to worsening symptoms
  \item Secondary infections or injuries
  \item Side effects of treatment
  \item Reduced quality of life
  \item Emotional and psychological impact
\end{itemize}
\section{Diagnosis}
To make the diagnosis, doctors look for tissue biopsy with Congo red staining and apple-green birefringence under polarized light---the diagnostic gold standard. Comprehensive medical history and physical examination Laboratory tests to assess organ function and rule out other conditions Imaging studies to visualize affected structures Specialized testing depending on the affected system Consultation with relevant specialists Monitoring of disease progression over time

\begin{itemize}[noitemsep]
  \item Comprehensive medical history and physical examination
  \item Laboratory tests to assess organ function
  \item Imaging studies to visualize affected structures
  \item Specialized testing depending on the affected system
  \item Consultation with relevant specialists
\end{itemize}
\section{Prevention}
\begin{itemize}[noitemsep]
  \item Maintain a healthy lifestyle with balanced diet and exercise
  \item Avoid known risk factors
  \item Attend regular medical check-ups for early detection
  \item Follow recommended screening guidelines
\end{itemize}
\section{Treatment Options}
Medications to manage symptoms and address underlying mechanisms Lifestyle modifications and self-care measures Physical therapy and rehabilitation Surgical intervention when indicated Regular monitoring and follow-up care Patient education and support services

\begin{itemize}[noitemsep]
  \item Medications to manage symptoms and address underlying mechanisms
  \item Lifestyle modifications and self-care measures
  \item Physical therapy and rehabilitation
  \item Surgical intervention when indicated
  \item Regular monitoring and follow-up care
\end{itemize}
\section{Living with It}
\begin{itemize}[noitemsep]
  \item Follow your treatment plan as prescribed
  \item Maintain regular follow-up with your healthcare provider
  \item Make necessary lifestyle adjustments
  \item Seek support from family, friends, or support groups
  \item Stay informed about your condition
\end{itemize}
\section{Outlook}
Unfortunately, the outlook is often poor for AL amyloidosis without treatment but has improved significantly with modern therapies. The outlook varies depending on the specific condition, severity at diagnosis, and response to treatment. Early intervention generally leads to better outcomes. Many conditions can be effectively managed with appropriate treatment. Regular follow-up care is important for monitoring and treatment adjustment.
